\documentclass[10pt]{beamer}
\usetheme[progressbar=frametitle]{metropolis}
\usepackage{appendixnumberbeamer}
\usepackage[style=verbose, backend=bibtex]{biblatex}
\addbibresource{slide.bib}

\usepackage{booktabs}
\usepackage[scale=2]{ccicons}

\usepackage{pgfplots}
\usepgfplotslibrary{dateplot}

\usepackage{xspace}
\newcommand{\themename}{\textbf{\textsc{metropolis}}\xspace}



\title{COMP0197: Applied Deep Learning - Group Project}
\subtitle{ProbTSLANet: Probabilistic Time Series Forecasting with Uncertainty Decomposition}
% \date{\today}
\author{Group 29} 
\institute{Chi Wang Cheung, Xuanhao Gu, Scarlett He, Baorui Li, Tongjia Li, Zhiyuan Luo, Shuhan Wang, \\ Flora Xu, Fengyuan Zheng}

\date{March 26, 2026}
\begin{document}

\maketitle

\begin{frame}{Table of contents}
  \setbeamertemplate{section in toc}[sections numbered]
  \tableofcontents%[hideallsubsections]
\end{frame}

\section[Design and Rationale]{Design and Rationale}

\begin{frame}[fragile]{Design and Rationale - Problem Definition}
\textbf{Project Goal}: 
Develop a sequence model predicting future values while quantifying uncertainties associated with those predictions.

\textbf{Motivation}: 
Standard forecasting models often only do point estimation, but lack a measure of \emph{confidence}. We aimed to 
build a system that outputs a \textbf{probability distribution} $\mathcal{N}(\mu, \sigma^2)$ for every future timestep.

\end{frame}

\begin{frame}[fragile]{Design and Rationale - Two types of uncertainty}

\begin{columns}[T,onlytextwidth]
    \footnote{{ Eyke Hüllermeier and Willem Waegeman. “Aleatoric and
    epistemic uncertainty in machine learning: An introduction to
    concepts and methods”. }}
    \column{0.5\textwidth}
      \textbf{Aleatoric Uncertainty}
      \begin{itemize}
        \item irreducible data noise 
        \item measurement error, chaos
        \newline
      \end{itemize}
      \centering
      \includegraphics[width=0.8\textwidth]{images/Aleatoric.png}
    
    \column{0.5\textwidth}
      \textbf{Epistemic Uncertainty}
      \begin{itemize}
        \item model uncertainty
        \item reducible with more data
        \newline
      \end{itemize}
      \centering
      \includegraphics[width=0.8\textwidth]{images/Epistemic.png}

  \end{columns}
   
\end{frame}

\begin{frame}[fragile]{Design and Rationale - Architecture Backbone} 

    Time Series Lightweight Adaptive Network \textbf{TSLANet} \textit{(ICML2024)} replaces Self-Attention with two lightweight, specialized blocks:

    \begin{columns}[T,onlytextwidth]
        \column{0.48\textwidth}
        \metroset{block=fill}
        \begin{block}{Adaptive Spectral Block (ASB)}
            \small
            \textbf{Frequency Domain Learning}
            \begin{itemize}
                \item Applies \textbf{FFT} to the input sequence.
                \item \textbf{Adaptive Masking}: Separates low-frequency trends from high-frequency noise based on energy.
                \item Applies \textbf{Inverse FFT} to retain it back to spatial domain
            \end{itemize}
        \end{block}

        \column{0.48\textwidth}
        \metroset{block=fill}
        \begin{block}{Interactive Conv Block (ICB)}
            \small
            \textbf{Time Domain Learning}
            \begin{itemize}
                \item Uses two parallel Convolutions (kernel sizes 1 and 3).
                \item \textbf{Gated Interaction}: Employs multiplicative gating to fuse features.
                \item Extracts local temporal dependencies and cross-variable correlations.
            \end{itemize}
        \end{block}
    \end{columns}
    
    \vspace{0.3cm}
    \centering
\end{frame}

\begin{frame}[fragile]{Design and Rationale - Architecture Backbone} 
    \begin{figure}
        \centering
        \includegraphics[width=0.85\textwidth]{images/TSLANet.png}
        \caption{\textit{Overall Architecture of TSLANet 
        \protect\footnote{Emadeldeen Eldele et al. “TSLAnet: Rethinking transformers for time series representation learning”. }}}
    \end{figure}
\end{frame}

\begin{frame}[fragile]{Design and Rationale - Probabilistic Extension}

\textbf{After ICB Backbone}

\[
\mathbf{h} = \text{Backbone}(x)
\]

We introduce two Gaussian output heads:
\[
\mu = W_{\mu}\mathbf{h} + b_{\mu}
\]
\[
\log \sigma^2 = W_{\sigma}\mathbf{h} + b_{\sigma}
\]

Heteroscedastic Gaussian Modeling with NLL Objective:

\[
\mathcal{L} = \frac{1}{2}\left(\log \sigma^2 + \frac{(y - \mu)^2}{\sigma^2}\right)
\]

\begin{itemize}
    \item $\mu$: predicted mean (forecast)
    \item $\sigma^2$: data-dependent uncertainty (aleatoric)
\end{itemize}

\end{frame}

\begin{frame}[fragile]{Design and Rationale - Epistemic Uncertainty via MC Dropout}

\textbf{Monte Carlo (MC) Dropout for Epistemic Uncertainty}

\vspace{0.5em}

At test time, we enable dropout and perform $K$ stochastic forward passes:
\[
\mu_k = f_{\theta, m_k}(x), \quad k = 1, \dots, K
\]

\vspace{0.5em}

\textbf{Epistemic uncertainty} (model uncertainty) is estimated as:
\[
\mathrm{Var}_{\text{ep}}(y \mid x) = \mathrm{Var}(\mu_1, \dots, \mu_K)
\]

\vspace{0.5em}

\textbf{Aleatoric uncertainty} (data noise) from Gaussian head:
\[
\mathrm{Var}_{\text{al}}(y \mid x) = \frac{1}{K} \sum_{k=1}^{K} \sigma_k^2
\]
\textbf{Total predictive uncertainty:}
\[
\mathrm{Var}_{\text{total}} = \mathrm{Var}_{\text{ep}} + \mathrm{Var}_{\text{al}}
\]

\end{frame}

\begin{frame}[fragile]{Design and Rationale - Dataset Choice} 

\begin{itemize}

\item \textbf{Dataset:} Weather benchmark — $\sim$52K hourly observations, 21 meteorological variables from a German weather station. Widely used in time-series forecasting literature (e.g., Autoformer, PatchTST).

\vspace{0.5em}

\item \textbf{Motivation:} Strong temporal dependency and clear periodic patterns (well-suited for sequence and spectral modelling), together with inherent stochastic noise (ideal for evaluating uncertainty quantification).

\vspace{0.5em}

\item \textbf{Data split:} Chronological 70/10/20 train/validation/test split to prevent information leakage.

\end{itemize}

\end{frame}


% ------------------------------------------------

\section{Implementation Plan}

\begin{frame}{Implementation Plan}
\small

\begin{columns}[T,onlytextwidth]

    % Left column
    \column{0.3\textwidth}
    \metroset{block=fill}   
    \begin{alertblock}{Framework}
    PyTorch + \texttt{pandas} + \texttt{matplotlib}
    \end{alertblock}

    \begin{exampleblock}{File Structure}
    \ttfamily
    \footnotesize
    src/\\
    \hspace*{1em}data/weather/\\
    \hspace*{1em}saved\_models/\\
    \hspace*{1em}run\_all.sh\\
    \hspace*{1em}test.py\\
    \hspace*{1em}train.py
    \end{exampleblock}
    
    % Right column
    \column{0.65\textwidth}
    \textbf{Implementation roadmap}
    \vspace{0.2cm}
    
    \begin{enumerate}
        \item \textbf{Data loader:} \texttt{Dataset\_Custom} for CSV input, z-score normalisation, sliding-window sampling, and time features
        
        \item \textbf{Model:} ProbabilisticTSLANet
        
        \item \textbf{Training loop:} self-supervised pretraining $\rightarrow$ supervised fine-tuning
        
        \item \textbf{Evaluation:} forecasting metrics and uncertainty plots
    \end{enumerate}

\end{columns}

\end{frame}

\begin{frame}{Implementation Plan - Evaluation}
\begin{columns}[T]
        \begin{column}{0.5\textwidth}
            \textbf{1. Point Metrics}
            \begin{itemize}
                \item \textbf{MSE}
                \item \textbf{MAE}
                \item \textbf{RMSE}
            \end{itemize}
            
            \textbf{2. Probabilistic Metrics}
            \begin{itemize}
                \item \textbf{CRPS}: Measures the integrated squared error between the predicted and empirical CDFs
                \item \textbf{NLL}: Measures the log-density of the observed data under the predicted distribution
            \end{itemize}
        \end{column}

        \begin{column}{0.5\textwidth}
            \textbf{3. Sharpness}
            \begin{itemize}
                \item \textbf{Sharpness}: Average width of the 90\%/50\% prediction intervals
            \end{itemize}

            \textbf{4. Calibration}
            \begin{itemize}
                \item \textbf{Calibration Error}: Mean absolute difference between predicted quantiles and actual coverage
            \end{itemize}
            

            \textbf{5. Uncertainty}
            \begin{itemize}
                \item \textbf{Mean Epistemic Var}
                \item \textbf{Mean Aleatoric Var}
            \end{itemize}
        \end{column}
    \end{columns}
\end{frame}

% ------------------------------------------------

\section{Current Progress}


\begin{frame}[shrink=8]{Current Progress — Experiment Design}
\begin{columns}[t]

  \begin{column}{0.49\textwidth}

    \begin{block}{\footnotesize\textbf{A1 — Deterministic Baseline}}
      \scriptsize
      ~\\
      \textbf{Principle:} Original TSLANet with no probabilistic modification.
      Outputs a \textbf{single point prediction} $\hat{y}$ per time step.
      No uncertainty is modelled. Serves as the reference for point accuracy.
      \par\vspace{4pt}
      \textbf{Method:} Trained with MSE loss. Single forward pass at inference;
      no distribution, no intervals.
    \end{block}

    \vspace{0.4em}

    \begin{block}{\footnotesize\textbf{A3 — Aleatoric Uncertainty (NLL Head)}}
      \scriptsize
      ~\\
      \textbf{Principle:} Models \textbf{aleatoric} uncertainty --- the
      inherent noise in the data itself, which cannot be reduced by more
      data or a better model. No MC Dropout; epistemic uncertainty is
      not captured.
      \par\vspace{4pt}
      \textbf{Method:} Output extended from $\hat{y}$ to $\mathcal{N}(\mu,\sigma^2)$.
      Training minimises NLL: $\mathcal{L} = \log\sigma + \frac{(y-\mu)^2}{2\sigma^2}$,
      directly supervising the model to learn how uncertain to be.
      Single forward pass at inference.
    \end{block}

  \end{column}

  \begin{column}{0.49\textwidth}

    \begin{block}{\footnotesize\textbf{A2 — Epistemic Uncertainty (MC Dropout)}}
      \scriptsize
      ~\\
      \textbf{Principle:} Models \textbf{epistemic} uncertainty --- the
      model's confidence in its own weights. More unstable outputs across
      passes indicate lower model confidence.
      \par\vspace{4pt}
      \textbf{Method:} Architecture and loss unchanged from A1.
      Dropout stays active at inference; $K$ forward passes each randomly
      silence a different subset of neurons. Mean of $K$ outputs gives
      the forecast; standard deviation estimates epistemic uncertainty.
    \end{block}

    \vspace{0.4em}

    \begin{block}{\footnotesize\textbf{A4 — Combined: NLL Head + MC Dropout}}
      \scriptsize
      ~\\
      \textbf{Principle:} Simultaneously models both \textbf{aleatoric}
      and \textbf{epistemic} uncertainty, aiming for a complete
      uncertainty budget.
      \par\vspace{4pt}
      \textbf{Method:} Combines A2 and A3. $K$ stochastic forward passes
      each produce $(\mu_k, \sigma^2_k)$; the two uncertainty sources
      are decomposed and summed to form the final predictive distribution.
    \end{block}

  \end{column}

\end{columns}
\end{frame}




\begin{frame}[t]{Current Progress — Prediction Intervals}

\vspace{0.3em}
\centering
\textbf{\footnotesize A1 — Deterministic Baseline: point prediction only, no intervals produced.}

\vspace{0.5em}

\begin{columns}[c]

  \begin{column}{0.32\textwidth}
    \centering
    \textbf{\footnotesize A2 — MC Dropout}\\[4pt]
    \includegraphics[width=\textwidth]{images/a2}\\[4pt]
    \scriptsize \textcolor{red!80!black}{Intervals static and collapsed.\\
    Width unchanged across horizon.\\
    Ground truth escapes entirely.}
  \end{column}

  \begin{column}{0.32\textwidth}
    \centering
    \textbf{\footnotesize A3 — NLL Head}\\[4pt]
    \includegraphics[width=\textwidth]{images/a3}\\[4pt]
    \scriptsize \textcolor{green!60!black}{Intervals fan out over horizon.\\
    Uncertainty grows with distance.\\
    Ground truth stays within 95\% PI.}
  \end{column}

  \begin{column}{0.32\textwidth}
    \centering
    \textbf{\footnotesize A4 — NLL + MC Dropout}\\[4pt]
    \includegraphics[width=\textwidth]{images/a4}\\[4pt]
    \scriptsize Same fan-out trend as A3.\\
    Intervals slightly wider overall;\\
    edges marginally more irregular.
  \end{column}

\end{columns}

\vspace{0.5em}
\centering
\textbf{\footnotesize $\Rightarrow$ MC Dropout + NLL did not achieve the expected effect.}\\[4pt]
\scriptsize
\textit{NLL and MSE optimise different objectives --- NLL trains the model to produce
a well-calibrated distribution by learning $\mu$ and $\sigma^2$, but does not directly
optimise point accuracy. Improving NLL does not necessarily improve MSE/MAE.}


\end{frame}


% ------------------------------------------------

\section[Problem Solving]{Problem Solving (technical hurdles)}

% ---- Metric comparison table ----
\begin{frame}[shrink=5]{Problem Solving — Key Metric Comparison}

\small
\textbf{Core issue:} NLL encourages the model to \textcolor{red}{\textbf{inflate variance}} to lower loss
instead of improving prediction accuracy ---
sacrificing MSE while producing \textcolor{red}{\textbf{meaningless}} uncertainty intervals.

\vspace{0.4em}

\begin{table}
\centering
\setlength{\tabcolsep}{6pt}
\begin{tabular}{lcccc}
\toprule
\textbf{Model} & \textbf{MSE} & \textbf{NLL} & \textbf{Cal.\ Error} & \textbf{Log Var} \\
\midrule
A1 (Deterministic)       & 0.1736 & ---    & ---    & ---    \\
A3 (NLL only)            & \textcolor{red}{\textbf{0.1968}} & 0.2861 & 0.1023 & moderate \\
A5 (NLL + aux.\ MSE)     & 0.1829 & 0.1674 & 0.1306 & moderate \\
A6 (Two-stage)           & 0.1730 & 0.1730 & 0.2562 & \textcolor{red}{\textbf{very large}} \\
A7 (A5 + MC Dropout)     & 0.1829 & 0.1357 & 0.1386 & moderate \\
\bottomrule
\end{tabular}
\end{table}

\vspace{0.3em}
\centering
\scriptsize
$\Rightarrow$ \textbf{A5 chosen as base}; extended to \textbf{A7} (A5 + MC Dropout)
for full aleatoric + epistemic uncertainty decomposition.

\end{frame}

% ---- A5 / A6 ----
\begin{frame}[shrink=8]{Problem Solving (technical hurdles)}
\begin{columns}[t]

  \begin{column}{0.49\textwidth}

    \begin{block}{\footnotesize\textbf{A5 — NLL + Auxiliary MSE Loss}}
      \scriptsize
      ~\\
      \textbf{Hurdle:} Pure NLL training encourages the model to
      \textbf{inflate} $\sigma^2$ to minimise loss, which
      \emph{dilutes the gradient} for $\mu$ and causes
      MSE to \textbf{rise} compared with the deterministic baseline~(A1).
      \par\vspace{4pt}
      \textbf{Solution:} Add a weighted MSE term alongside NLL:
      \[
        \mathcal{L} = \mathcal{L}_{\text{NLL}} + \lambda\,\mathcal{L}_{\text{MSE}}
      \]
      This explicitly penalises poor mean predictions,
      preventing variance from masking inaccuracy while
      preserving well-calibrated uncertainty intervals.
    \end{block}

  \end{column}

  \begin{column}{0.49\textwidth}

    \begin{block}{\footnotesize\textbf{A6 — Two-Stage Decoupled Training}}
      \scriptsize
      ~\\
      \textbf{Hurdle:} Even with auxiliary MSE, jointly optimising
      $\mu$ and $\sigma^2$ allows the variance head to
      \textbf{compensate} for a weak mean by growing $\log\sigma^2$,
      producing intervals too wide to be informative.
      \par\vspace{4pt}
      \textbf{Solution:} Decouple training into two stages:
      \begin{enumerate}
        \item Train backbone + mean head with \textbf{MSE only};
              weights \textbf{frozen} after convergence.
        \item Train variance head alone with \textbf{NLL only},
              forcing it to learn \emph{residual uncertainty}
              rather than compensating for mean error.
      \end{enumerate}
      \textbf{Outcome:} Mean accuracy preserved; however variance
      head still tends toward overly large $\log\sigma^2$
      $\Rightarrow$ A5 selected as stronger base for A7.
    \end{block}

  \end{column}

\end{columns}
\end{frame}


% ------------------------------------------------




\section{Conclusion}

\begin{frame}[fragile]{Summary} 

  \begin{itemize}
      \item \textbf{ProbTSLANet} extends TSLANet with \textbf{Gaussian output heads} and \textbf{MC Dropout}
  \end{itemize}
\begin{itemize}
    \item Add \textbf{auxiliary MSE loss} on top of NLL loss, improving MSE while maintaining calibration

\end{itemize}
\begin{itemize}

    \item Our model produces calibrated probabilistic weather forecasts with decomposed aleatoric/epistemic uncertainty
\end{itemize}



\end{frame}


% ------------------------------------------------




\section{Thank you for listening!}








% ------------------------------------------------


\section{Elements}

\begin{frame}[fragile]{Typography}
      \begin{verbatim}The theme provides sensible defaults to
\emph{emphasize} text, \alert{accent} parts
or show \textbf{bold} results.\end{verbatim}

  \begin{center}becomes\end{center}

  The theme provides sensible defaults to \emph{emphasize} text,
  \alert{accent} parts or show \textbf{bold} results.
\end{frame}

\begin{frame}{Font feature test}
  \begin{itemize}
    \item Regular
    \item \textit{Italic}
    \item \textsc{SmallCaps}
    \item \textbf{Bold}
    \item \textbf{\textit{Bold Italic}}
    \item \textbf{\textsc{Bold SmallCaps}}
    \item \texttt{Monospace}
    \item \texttt{\textit{Monospace Italic}}
    \item \texttt{\textbf{Monospace Bold}}
    \item \texttt{\textbf{\textit{Monospace Bold Italic}}}
  \end{itemize}
\end{frame}

\begin{frame}{Lists}
  \begin{columns}[T,onlytextwidth]
    \column{0.33\textwidth}
      Items
      \begin{itemize}
        \item Milk \item Eggs \item Potatos
      \end{itemize}

    \column{0.33\textwidth}
      Enumerations
      \begin{enumerate}
        \item First, \item Second and \item Last.
      \end{enumerate}

    \column{0.33\textwidth}
      Descriptions
      \begin{description}
        \item[PowerPoint] Meeh. \item[Beamer] Yeeeha.
      \end{description}
  \end{columns}
\end{frame}
\begin{frame}{Animation}
  \begin{itemize}[<+- | alert@+>]
    \item \alert<4>{This is\only<4>{ really} important}
    \item Now this
    \item And now this
  \end{itemize}
\end{frame}
\begin{frame}{Figures}
  \begin{figure}
    \newcounter{density}
    \setcounter{density}{20}
    \begin{tikzpicture}
      \def\couleur{alerted text.fg}
      \path[coordinate] (0,0)  coordinate(A)
                  ++( 90:5cm) coordinate(B)
                  ++(0:5cm) coordinate(C)
                  ++(-90:5cm) coordinate(D);
      \draw[fill=\couleur!\thedensity] (A) -- (B) -- (C) --(D) -- cycle;
      \foreach \x in {1,...,40}{%
          \pgfmathsetcounter{density}{\thedensity+20}
          \setcounter{density}{\thedensity}
          \path[coordinate] coordinate(X) at (A){};
          \path[coordinate] (A) -- (B) coordinate[pos=.10](A)
                              -- (C) coordinate[pos=.10](B)
                              -- (D) coordinate[pos=.10](C)
                              -- (X) coordinate[pos=.10](D);
          \draw[fill=\couleur!\thedensity] (A)--(B)--(C)-- (D) -- cycle;
      }
    \end{tikzpicture}
    \caption{Rotated square from
    \href{http://www.texample.net/tikz/examples/rotated-polygons/}{texample.net}.}
  \end{figure}
\end{frame}
\begin{frame}{Tables}
  \begin{table}
    \caption{Largest cities in the world (source: Wikipedia)}
    \begin{tabular}{lr}
      \toprule
      City & Population\\
      \midrule
      Mexico City & 20,116,842\\
      Shanghai & 19,210,000\\
      Peking & 15,796,450\\
      Istanbul & 14,160,467\\
      \bottomrule
    \end{tabular}
  \end{table}
\end{frame}
\begin{frame}{Blocks}
  Three different block environments are pre-defined and may be styled with an
  optional background color.

  \begin{columns}[T,onlytextwidth]
    \column{0.5\textwidth}
      \begin{block}{Default}
        Block content.
      \end{block}

      \begin{alertblock}{Alert}
        Block content.
      \end{alertblock}

      \begin{exampleblock}{Example}
        Block content.
      \end{exampleblock}

    \column{0.5\textwidth}

      \metroset{block=fill}

      \begin{block}{Default}
        Block content.
      \end{block}

      \begin{alertblock}{Alert}
        Block content.
      \end{alertblock}

      \begin{exampleblock}{Example}
        Block content.
      \end{exampleblock}

  \end{columns}
\end{frame}
\begin{frame}{Math}
  \begin{equation*}
    e = \lim_{n\to \infty} \left(1 + \frac{1}{n}\right)^n
  \end{equation*}
\end{frame}
\begin{frame}{Line plots}
  \begin{figure}
    \begin{tikzpicture}
      \begin{axis}[
        mlineplot,
        width=0.9\textwidth,
        height=6cm,
      ]

        \addplot {sin(deg(x))};
        \addplot+[samples=100] {sin(deg(2*x))};

      \end{axis}
    \end{tikzpicture}
  \end{figure}
\end{frame}
\begin{frame}{Bar charts}
  \begin{figure}
    \begin{tikzpicture}
      \begin{axis}[
        mbarplot,
        xlabel={Foo},
        ylabel={Bar},
        width=0.9\textwidth,
        height=6cm,
      ]

      \addplot plot coordinates {(1, 20) (2, 25) (3, 22.4) (4, 12.4)};
      \addplot plot coordinates {(1, 18) (2, 24) (3, 23.5) (4, 13.2)};
      \addplot plot coordinates {(1, 10) (2, 19) (3, 25) (4, 15.2)};

      \legend{lorem, ipsum, dolor}

      \end{axis}
    \end{tikzpicture}
  \end{figure}
\end{frame}
\begin{frame}{Quotes}
  \begin{quote}
    Veni, Vidi, Vici
  \end{quote}
\end{frame}

{%
\setbeamertemplate{frame footer}{My custom footer}
\begin{frame}[fragile]{Frame footer}
    \themename defines a custom beamer template to add a text to the footer. It can be set via
    \begin{verbatim}\setbeamertemplate{frame footer}{My custom footer}\end{verbatim}
\end{frame}
}

\begin{frame}{References}
  Some references to showcase [allowframebreaks] \cite{knuth92,ConcreteMath,Simpson,Er01,greenwade93}
\end{frame}

\section{Conclusion}

\begin{frame}[fragile]{Summary} 

  \begin{itemize}

\end{itemize}



\end{frame}

{\setbeamercolor{palette primary}{fg=black, bg=yellow}
\begin{frame}[standout]
  Questions?
\end{frame}
}

\appendix

\begin{frame}[fragile]{Backup slides}
  Sometimes, it is useful to add slides at the end of your presentation to
  refer to during audience questions.

  The best way to do this is to include the \verb|appendixnumberbeamer|
  package in your preamble and call \verb|\appendix| before your backup slides.

  \themename will automatically turn off slide numbering and progress bars for
  slides in the appendix.
\end{frame}

\begin{frame}[allowframebreaks]{References}
  \footnotesize
  \bibliographystyle{abbrv}
  \bibliography{slide}
\end{frame}

\end{document}
